\chapter{العصبونات بحسب عدد التغصنات}
\chaptermark{عدد التغصنات}
\label{sec:dendrites}

\section{عدد المداخل معاملًا في الدارة}

صنّفنا العصبونات في الفصل~\ref{sec:neurontypes} بحسب ما يطلقه مخرجها، وفي الفصل~\ref{sec:eetypes} بحسب الجهاز الذي تعمل عمله. وثمة معيار ثالث يلاحظه المهندس على الفور: \emph{كم مدخلًا للعصبون}. تجمع المداخلَ التغصناتُ، وهي الفروع التي تحط عليها محاوير الخلايا الأخرى (الفصل~\ref{sec:dofa}، القسم~\ref{sec:weightstore}). بعض العصبونات بلا تغصنات أصلًا، وبعضها له تغصن أو اثنان، وبعضها له شجرة تجمع مئة ألف مدخل. هذا عند مهندس الإلكترونيات عامل التحميل عند المدخل، أي \foreignlanguage{english}{fan-in}، وهو يكاد يوافق المهمة دائمًا.

لماذا يهم عدد التغصنات وحجمها إلى هذا الحد؟ يتضح ذلك من تقدير بسيط. الغشاء مكثف سعته النوعية $c_m\approx1\,\text{\textmu F/cm}^2$، وكلما كبرت مساحة العصبون $A$ مع تغصناته، كبرت الشحنة التي يجب جلبها لرفع الجهد بمقدار $\Delta V$. والزمن الذي يفعل فيه تيار الدخل $I$ ذلك، إن أهملنا التسريب، هو
\begin{empheq}[box=\keybox]{equation}
t \;\approx\; \frac{c_m\,A\,\Delta V}{I} .
\label{eq:dendriteload}
\end{empheq}
العصبون الصغير ذو التغصنات القصيرة القليلة مساحته نحو $2\cdot10^{-5}\,\text{cm}^2$ وسعته نحو $20\,\text{pF}$. ومشبك واحد كبير بتيار بضعة نانوأمبيرات يرفعه $20\mV$ في أقل من عُشر جزء من الألف من الثانية. أما العصبون ذو الشجرة الكبيرة فمساحته أكبر بنحو خمس عشرة مرة وسعته نحو $300\,\text{pF}$؛ والمشبك العادي بتيار نحو $0.1\,\text{nA}$ كان سيشحنه في $60\ms$، وهو أطول بكثير من الثابت الزمني للتسريب، أي أنه لن يشحنه أبدًا. مثل هذا العصبون يحتاج إلى عشرات المداخل المتزامنة. الأعداد هنا رتب مقادير، لكن النتيجة لا تعتمد عليها: \emph{مداخل قليلة، سرعة ودقة؛ مداخل كثيرة، بطء وعمل بالمجموع}.

\begin{figure}[t]
\centering
\begin{tikzpicture}[>=Stealth,
  soma/.style={circle,draw,very thick,fill=blue!6,minimum size=0.55cm,inner sep=0pt},
  ax/.style={->,very thick,red!70!black},
  den/.style={very thick,blue!60!black}]
% 0 dendrites: sensor on a line
\begin{scope}[xshift=0cm]
\draw[den,decorate,decoration={snake,amplitude=1pt,segment length=4pt}] (-0.9,0) -- (0,0);
\node[font=\scriptsize] at (-0.9,0.3) {مستشعر};
\draw[ax] (0,0) -- (1.0,0);
\draw[thick] (0.1,0) -- (0.1,0.45);
\node[soma,minimum size=0.4cm] at (0.1,0.65) {};
\node[font=\scriptsize,align=center] at (0.05,-0.75) {$0$\\خط بمستشعر};
\end{scope}
% 1 dendrite
\begin{scope}[xshift=3.3cm]
\draw[den] (-0.9,0) -- (-0.28,0);
\node[soma] at (0,0) {};
\draw[ax] (0.28,0) -- (0.9,0);
\node[font=\scriptsize,align=center] at (0,-0.75) {$1$\\عازل، مكرِّر};
\end{scope}
% 2 dendrites: one per ear
\begin{scope}[xshift=6.2cm]
\draw[den] (-0.28,0) -- (-0.9,0); \draw[den] (0.28,0) -- (0.9,0);
\node[soma] at (0,0) {};
\draw[ax] (0,-0.28) -- (0,-0.6);
\node[font=\scriptsize] at (-0.9,0.25) {يسار}; \node[font=\scriptsize] at (0.9,0.25) {يمين};
\node[font=\scriptsize,align=center] at (0,-1.0) {$2$\\AND زمنية};
\end{scope}
% ~4 short dendrites: mixer
\begin{scope}[xshift=9.0cm]
\foreach \a in {45,135,225,315}{\draw[den] (\a:0.28) -- (\a:0.7); \fill[blue!60!black] (\a:0.72) circle (0.06);}
\node[soma] at (0,0) {};
\draw[ax] (0.28,0) -- (1.0,0);
\node[font=\scriptsize,align=center] at (0,-1.0) {$\sim4$\\مازج};
\end{scope}
% many: summing tree
\begin{scope}[xshift=12.0cm]
\foreach \a in {60,90,120}{\draw[den] (0,0.28) -- ++(\a:0.55) coordinate (b); \foreach \d in {-20,20}{\draw[den] (b) -- ++(\a+\d:0.35);}}
\foreach \a in {-120,-60}{\draw[den] (0,-0.28) -- ++(\a:0.45);}
\node[soma] at (0,0) {};
\draw[ax] (0.28,0) -- (1.0,0);
\node[font=\scriptsize,align=center] at (0,-1.0) {$10^4\text{--}10^5$\\جامع};
\end{scope}
\end{tikzpicture}
\caption{خمسة أصناف بحسب عدد المداخل. بالأزرق التغصنات (المداخل)، وبالأحمر المحوار (المخرج). كلما قلّت المداخل، نقل العصبون الإشارة المفردة أسرع وأدق؛ وكلما كثرت، أحسن الجمع والتصنيف. هذا مخطط، لا رسم لأشكال الخلايا.}
\label{fig:dendriteclasses}
\end{figure}

\section{صفر تغصنات: مستشعر في طرف الخط}

العصبونات الحسية للمس والألم وتمدد العضلات (الفصول~\ref{sec:sensors} و\ref{sec:pain} و\ref{sec:muscles}) ليس على جسم خليتها تغصن واحد. يخرج المحوار من الجسم فرعًا واحدًا وينقسم فورًا إلى اثنين: فرع يذهب إلى الجلد أو العضلة، وطرفه نفسه هو المستشعر؛ وفرع آخر يذهب إلى الحبل الشوكي. تجري النبضة من المستشعر إلى الحبل الشوكي مباشرة، متجاوزة جسم الخلية. هذا عند المهندس خط اتصال مستشعرُه في الطرف البعيد، وجسم الخلية معلّق جانبًا كوحدة تغذية: يمدّ الخط بالطاقة، لكنه لا يشارك في النقل. والمكسب هو السرعة والموثوقية: لا يوجد على الطريق موضع واحد يلزم فيه جمع الإشارة وتقرير إرسالها من جديد.

ومستشعرات الضوء والصوت حالة أكثر تطرفًا: فهي ليست عصبونات جامعة، بل محوّلات بحد ذاتها، مدخلها ليس مشبكًا بل بروتين مستقبِل (الشكل~\ref{fig:transducer}).

\section{تغصن واحد: عازل}

العصبون ذو التغصن الواحد والمحوار الواحد يتلقى خط دخل واحدًا وينقله إلى الأمام. هكذا بُنيت عصبونات الشم: تغصنها الوحيد يخرج إلى سطح الأنف ويحمل نوعًا واحدًا من المستقبلات~\cite{Mombaerts2004}؛ وهكذا بُنيت خلايا النقل في الشبكية بين مستشعرات الضوء والعصبونات المخرجة للعين؛ وهكذا بُنيت العصبونات التي تصل مستشعرات الأذن بالدماغ. كان المهندس سيسميها عازلًا أو مكرِّرًا: فهي لا تحسب، بل توصل إشارة واحدة، وقد تغيّر شكلها أحيانًا، كأن تحوّل الجهد المتصل للمستشعر إلى نبضات، كما في الفصل~\ref{sec:sensors}.

\section{تغصنان: بوابة AND زمنية}

كواشف التزامن التي تحدد اتجاه الصوت (الشكل~\ref{fig:itd}) كثيرًا ما يكون لها عند الثدييات تغصنان بالضبط، متجهان في اتجاهين متعاكسين: إلى أحدهما تصل المداخل من الأذن اليسرى، وإلى الآخر من اليمنى~\cite{Grothe2010}. لماذا تُفرَّق المداخل؟ لنتذكر الصيغة~\eqref{eq:shunt}: حين تنفتح قنوات مثيرة كثيرة في موضع واحد من الغشاء، يقترب الجهد هناك من جهد التوازن، ويضيف كل مدخل تالٍ أقل فأقل. لذلك فالمداخل الكثيرة من أذن واحدة، المجموعة على تغصن واحد، تُشبعه ولا تستطيع وحدها بلوغ العصبون العتبة. أما مدخل واحد من كل تغصن فيُجمعان خطيًا تقريبًا. التغصنان يحوّلان العصبون إلى بوابة AND أكثر صرامة: أطلِق فقط إذا جاءت الإشارتان من الجانبين معًا. وقد بيّنت النماذج أن هذا التفريق يحسّن فعلًا كشف التزامن~\cite{AgmonSnir1998}.

\section{تغصنات قصيرة قليلة: المازج ومُعيد الإرسال}

\emph{المازجات.} في دارة التعلم الحركي في الفصل~\ref{sec:motorlearning} يوسّع نحو سبعين مليار عصبون \nt{GLUT} صغير المدخلَ إلى متجه طويل جدًا. لكل منها نحو أربعة تغصنات قصيرة فقط، وينتهي كل تغصن بـ”مخلب“ يطوّق نهاية مدخل واحد. وهكذا لا يرى كل مازج إلا نحو أربعة مداخل من مداخل كثيرة، ويعمل بوابة AND صغيرة على رباعيته. من $M$ خط دخل يمكن اختيار $\binom{M}{4}$ رباعية مختلفة: عند $M=100$ يقارب ذلك أربعة ملايين. لماذا هذا العدد كله؟ العنصر العتبي ذو $n$ مدخلًا يستطيع أن يفصل فصلًا صحيحًا إلى صنفين ما لا يزيد على
\begin{empheq}[box=\keybox]{equation}
P_{\max} \;\approx\; 2n
\label{eq:cover}
\end{empheq}
نمطًا عشوائيًا~\cite{Cover1965}. العصبون المخرج في الدارة نفسها يتلقى نحو $10^5$ مدخل من المازجات، والحد الأعلى له وفق~\eqref{eq:cover} من رتبة $2\cdot10^5$ تقابل مختلف. هذا حد العنصر المثالي: إذا كانت كل الأوزان موجبة، كما في المشابك المثيرة، ولزم هامش للضجيج، فالتقدير للعصبون نفسه نحو $4\cdot10^4$ تقابل~\cite{Brunel2004}. المازجات ذات المداخل القليلة لازمة تحديدًا كي يكون للجامع الكبير مداخل كثيرة مختلفة شبه مستقلة~\cite{Marr1969,Albus1971}.

\emph{مُعيدات الإرسال.} على الطريق من الأذن إلى كواشف الاتجاه عصبونات بتغصنات قصيرة قليلة، تتلقى بضعة مشابك ضخمة فقط، وبعضها في الجوهر مشبكًا واحدًا. ووفق~\eqref{eq:dendriteload} هذا بالضبط ما يلزم: سعة صغيرة وتيار كبير، فيطلق العصبون بعد أجزاء من جزء الألف من الثانية من كل نبضة داخلة، دون أن يفقد تقريبًا شيئًا من دقة التوقيت. أحد مُعيدات الإرسال هذه يتلقى \nt{GLUT} ويُخرج \nt{GLYC}: إنه عاكس بدقة عشرات الميكروثواني، يزوّد كواشف الاتجاه بتثبيط سريع~\cite{BorstSoria2012}.

\section{آلاف المداخل: الجامع والمصنِّف}

في الطرف الآخر من السلّم عصبونات \nt{GLUT} القشرية ذات العشرة آلاف مدخل، وعصبونات \nt{GABA} المخرجة في دارة التعلم الحركي ذات المئة ألف. هذا العصبون بطيء وفق~\eqref{eq:dendriteload} ولا يستطيع الاستجابة لمدخل مفرد: إنه يجمع. إنه العصبون المكامل في الفصل~\ref{sec:glutcomputer}، والجامع الذي تُبنى منه المصنِّفات. والشجرة الكبيرة تضيف شيئًا جديدًا أيضًا: تعمل فروعها دارات فرعية مستقلة لكل منها عتبتها، فيتصرف العصبون الواحد كشبكة صغيرة متعددة الطبقات~\cite{Beniaguev2021,Gidon2020}.

\section{تغصنات هي المخرج أيضًا}

ثمة عصبونات لا محوار لها أصلًا، وتغصناتها هي المدخل والمخرج معًا: تتلقى الإشارات وتطلق الناقل العصبي بنفسها. وأكثر الأمثلة دراسةً عصبونات \nt{GABA} في الشبكية المشاركة في تحديد اتجاه الحركة (الفصل~\ref{sec:gaba}، الشكل~\ref{fig:direction}). في هذا العصبون يستجيب كل فرع بمفرده للحركة من جسم الخلية نحو طرفه، ويُصدر التثبيط من هذا الطرف~\cite{Euler2002}. العصبون الواحد عدة كواشف اتجاه مستقلة، واحد لكل فرع. وهو عند المهندس ليس عنصرًا واحدًا، بل دارة متكاملة بعدة قنوات في غلاف واحد.

\section{الخلاصة}

\begin{table}[t]
\centering
\footnotesize
\setlength{\tabcolsep}{3pt}
\renewcommand{\arraystretch}{1.15}
\begin{tabular}{>{\raggedright}p{1.9cm}>{\raggedright}p{3.4cm}>{\raggedright}p{2.6cm}>{\raggedright}p{1.8cm}>{\raggedright\arraybackslash}p{3.8cm}}
\toprule
التغصنات & مثال & الجهاز & المداخل & ما المعروف \\
\midrule
$0$ & مستشعرات اللمس والألم والتمدد & خط بمستشعر في طرفه & --- & ثابت \\
$1$ & عصبونات الشم، خلايا النقل في الشبكية، عصبونات الأذن & عازل، مكرِّر & من $1$ إلى بضعة & ثابت \\
$2$ & كواشف اتجاه الصوت & بوابة AND زمنية & حزمتان & البنية ثابتة؛ ومكسب تفريق المداخل مبيَّن في النماذج \\
$\sim4$ & مازجات دارة التعلم الحركي & بوابة AND صغيرة & $\sim4$ & ثابت؛ ودور توسيع المدخل فرضية جيدة الأساس \\
قليلة، قصيرة & مُعيدات الإرسال على الطريق من الأذن & مُعيد إرسال، عاكس & من $1$ إلى بضعة ضخمة & ثابت \\
آلاف & عصبونات \nt{GLUT} القشرية، العصبونات المخرجة للتعلم الحركي & جامع، مصنِّف & $10^4\text{--}10^5$ & ثابت؛ والفروع دارات فرعية مُقاسة في أمثلة منفردة \\
تغصنات مُخرِجة & عصبونات \nt{GABA} للاتجاه في الشبكية & دارة متكاملة متعددة القنوات & كثيرة & مُقاس \\
\bottomrule
\end{tabular}
\caption{العصبونات بحسب عدد التغصنات.}
\label{tab:dendrites}
\end{table}

\begin{keyblock}
\begin{proposition}[عدد المداخل يتبع المهمة]
\label{prop:dendrites}
حيث تلزم سرعة الإشارة المفردة ودقتها، في المستشعرات ومُعيدات الإرسال وكواشف التزامن، يكون للعصبونات تغصنات قليلة ومداخل قليلة ومشابك كبيرة: السعة الصغيرة~\eqref{eq:dendriteload} تُشحن بسرعة. وحيث يلزم الجمع والتصنيف، يكون للعصبونات أشجار كبيرة بآلاف المداخل. بنية الأصناف مُقاسة؛ والتفسير الحسابي جيد الأساس، لكنه يبقى فرضية لأصناف كثيرة.
\end{proposition}
\end{keyblock}

يكمّل الجدول~\ref{tab:dendrites} التصنيفين السابقين: بحسب الناقل العصبي (الجدول~\ref{tab:types}) وبحسب الجهاز (الجدول~\ref{tab:eetypes}). والثلاثة تنظر إلى العنصر نفسه من جهات مختلفة: ماذا يُخرج، وماذا يحسب، وإلى كم يستمع.

\section{تمارين}

\begin{exercise}[\lvlA]
\label{ex:19:1}
\emph{كم مدخلًا يحتاج العصبون الكبير.} عصبون بـ$C=300\,\text{pF}$ و$\tau_m=20\ms$ يتلقى مشابك هي مصادر تيار كل منها $0.1\,\text{nA}$، والعتبة أعلى من الراحة بـ$20\mV$. (أ)~كم مشبكًا يجب أن يعمل معًا لبلوغ العتبة أصلًا؛ وفي $10\ms$؛ وفي $5\ms$؟ (ب)~في كم يبلغ العتبةَ عصبونٌ بـ$C=20\,\text{pF}$ من مشبك $4\,\text{nA}$؟
\end{exercise}

\begin{exercise}[\lvlA]
\label{ex:19:2}
\emph{لماذا أربعة.} المازج بوابة AND على $k$ خطًا من $M=100$، ونصف الخطوط نشط في النمط، والمازجات $7\cdot10^{10}$. (أ)~كم منها يستجيب للنمط عند $k=2$؛ $4$؛ $40$؟ (ب)~كم ضعفًا تزيد المازجات عدد التقابلات~\eqref{eq:cover} للعصبون المخرج ذي $10^5$ مدخل مقارنة بـ$M=100$ خط مباشر؟
\end{exercise}

\begin{exercise}[\lvlB]
\label{ex:19:3}
\emph{من أين يأتي $2n$.} المستوى الفائق المار بالأصل يقسم $P$ نقطة في وضع عام في فضاء ذي $n$ بعدًا إلى صنفين بـ$C(P,n)=2\sum_{k=0}^{n-1}\binom{P-1}{k}$ طريقة. (أ)~أثبت أنه عند $P=2n$ يكون نصف التقسيمات الـ$2^P$ بالضبط قابلًا للفصل. (ب)~ما النسبة القابلة للفصل عند $n=100$ و$P=180$؛ $220$؟
\end{exercise}

\begin{exercise}[\lvlB]
\label{ex:19:4}
\emph{التغصنان بوابةَ AND صارمة.} التغصن حجيرة بتسريب $g_L$، وكل مدخل يفتح عليه موصلية $g_0=0.5\,g_L$ بجهد $E_E$، والجسم يرى $\kappa(V_1+V_2)$، والعتبة $\theta=1.1\,\kappa E_E$. لماذا لا يطلق العصبونَ أيُّ عدد من المداخل على تغصن واحد، وكم مدخلًا يلزم على كل تغصن؟
\end{exercise}

\begin{exercise}[\lvlB]
\label{ex:19:5}
\emph{تصميم مُعيد إرسال.} ارتعاش النبضة $\sigma_t=\sigma_V/\dot V$، حيث $\dot V$ ميل الجهد عند العتبة. المطلوب $\sigma_t\le20\,\text{\textmu s}$ عند ضجيج $\sigma_V=1\mV$؛ أهمل التسريب. كم مشبكًا متزامنًا بـ$0.1\,\text{nA}$ يلزم لذلك عصبونًا بـ$C=300\,\text{pF}$، وما ارتعاش عصبون بـ$C=20\,\text{pF}$ ومشبك واحد $4\,\text{nA}$؟
\end{exercise}

\begin{exercise}[\lvlB]
\label{ex:19:6}
\emph{نمذجة: شحن تغصن طويل.} حاكِ كابلًا منفعلًا $\tau_m\partial_tV=\lambda^2\partial_x^2V-V+i/g$ طوله $L=\ell/\lambda$ بطرفين مغلقين ($40$ حجيرة، أويلر). تُحقن في الطرف البعيد نبضة تيار قصيرة: متى، وإلى أي نسبة من الجهد الناتج عن الشحنة نفسها موزعة على التغصن كله، يرتفع الطرف القريب عند $L=0.2$؛ $1$؛ $2$؟
\end{exercise}

\begin{exercise}[\lvlC]
\label{ex:19:7}
\emph{بحث: العدد الأمثل للمداخل.} السعة $C=C_0+nc_s$، ويعمل معًا $m$ مدخلًا تيار كل منها $I_s$، ويحفظ العصبون $P\approx2n$ تقابلًا~\eqref{eq:cover}. كيف يعتمد زمن الاستجابة على $P$ عند $m$ ثابت وعند نسبة ثابتة $m=\varphi n$؟ اقترح معيار كلفة وأوجد $n$ الأمثل لمُعيد الإرسال وللمصنِّف.
\end{exercise}
